
We designed a validation study to measure data availability across model families, benchmarks, and timepoints.

\textbf{Model selection.} We selected three model families from independent organizations to test whether category-level reporting is systematic across the LLM ecosystem:
\begin{itemize}
\item GPT family (OpenAI): GPT-3.5-turbo (baseline, 2022) and GPT-4 (current, 2023)
\item Claude family (Anthropic): Claude-2 (baseline, 2023) and Claude-3 (current, 2024)
\item Llama family (Meta): Llama-2 (baseline, 2023) and Llama-3 (current, 2024)
\end{itemize}

These organizations differ in institutional structure (for-profit with closed models, public benefit corporation, open-source-oriented lab) and represent independent development efforts.

\textbf{Benchmark selection.} We selected TruthfulQA and MMLU as two widely-reported benchmarks with explicit category structure in their design. TruthfulQA contains 817 questions across 38 original categories, which reports typically aggregate. MMLU contains 14,042 questions across 57 subjects.

\textbf{Temporal coverage.} For each family, we extracted data for baseline models (released 2022-2023) and current models (released 2023-2024), providing approximately 1-2 years separation for potential longitudinal analysis.

\textbf{Extraction procedure.} We downloaded technical reports from official sources (OpenAI website, Anthropic documentation, Meta research portal). We manually extracted category-level performance data from tables in these reports. Where reports provided accuracy, we converted to error rate (1 - accuracy). We normalized schema variations across reports (different column names, table structures) to a common format. We cross-referenced extracted values with report text to verify accuracy.

We used manual extraction rather than automated PDF parsing due to known limitations of parsing libraries on complex table structures. This represents a methodological choice prioritizing accuracy over scalability. For a data availability study, confirming that data exist is the primary requirement; automated extraction is an engineering challenge that could be addressed in future work focused on scalability.

\textbf{Validation criteria.} We defined four quantitative gates:
\begin{itemize}
\item Gate 1 (Family coverage): ≥3 model families with category-level data for both timepoints
\item Gate 2 (Granularity): ≥10 categories per benchmark
\item Gate 3 (Completeness): ≥90\% of expected data cells present
\item Gate 4 (Temporal consistency): Both baseline and current data available for all families
\end{itemize}

These thresholds were established before data extraction based on requirements from weak supervision literature (minimum 8-12 categories for robust clustering) and statistical analysis practices (tolerance for some missing data while ensuring robustness).
